\chapter{Ulm: metoden som kollapsa}

\begin{motto}
Det er lærerikt at det var her faget kom nærast eit fundament.\\ Det er endå meir lærerikt at det var her det rasa.
\end{motto}

Av alle stadene i denne historia er det ein som fortener eit særskilt vakt blikk, fordi det var her designfaget kom nærare eit fagleg fundament enn nokon gong før eller sidan. Hochschule für Gestaltung i Ulm, opna i 1953 og stengd i 1968, sette seg eit mål dei andre rørslene berre hadde antyda: å gjere formgjeving til ein disiplin med metode, med teori, med målbart grunnlag — ein praksis ein kunne grunngje og prøve, ikkje berre utøve. At dette forsøket vart det mest seriøse i fagets historie, og at det likevel kollapsa, både institusjonelt og fagleg, er ikkje ein tilfeldig tragedie. Det er ei av dei klåraste leksene boka har å gje. For når det beste forsøket på eit fundament både kjem nærast og fell, må ein spørje om det var noko ved sjølve oppgåva, og ikkje berre ved omstenda, som ikkje heldt.

\section{Frå Bauhaus til vitskap}

Ulm vart unnfanga som ein arv og vart noko anna. Skulen voks ut av etterkrigstidas tyske gjenreising, grunnlagd dels til minne om motstandsfolk avretta av nazistane, og dei fyrste åra stod under leiing av Max Bill, sjølv ein Bauhaus-elev. Bill ville at Ulm skulle vere eit nytt Bauhaus: kunst og handverk og industri sameina, formgjevaren som ein kunstnar med samfunnsansvar. Men ein yngre generasjon, med argentinaren Tomás Maldonado i spissen, ville lenger. Dei meinte at Bauhaus-arven, med sin tru på kunstnaren sin intuisjon, var utilstrekkeleg for ein industri som var vorten umåteleg meir kompleks. Maldonado sette det reint ut: opplæringa av formgjevarar måtte byggjast på ny kunnskap frå industri og vitskap, ikkje på det kunstnariske geniet sin teft\autocite{maldonado1958}. Striden enda med at Bill forlét skulen i 1957, og med det vende Ulm ryggen til kunsten og andletet mot vitskapen.

Det som no kom inn i timeplanen, var det ingen designskule hadde undervist før. Studentane las semiotikk og lærte å analysere teikn og tyding. Dei las ergonomi og menneskelege faktorar, og målte kroppen mot maskinen. Dei las informasjonsteori, kybernetikk, systemteori, produksjonsøkonomi, vitskapsteori. Maldonado og kollegaene hans ville gje formgjevaren eit grunnlag som kunne forsvarast med same strengskap som ein ingeniør sitt — ei \emph{vitskapleg operasjonalisering} av designprosessen, der vala ikkje lenger sprang ut av smak, men ut av analyse. Herbert Lindinger, som sjølv var elev og lærar der, har dokumentert kor systematisk dette vart bygt ut: kvart fag på skulen skulle tene ein metode, og metoden skulle gjere formgjevinga etterprøvbar\autocite{lindinger1991ulm}. Det var, så langt det gjekk, det mest gjennomtenkte forsøket faget har gjort på å bli noko ein kan kalle ein vitskap.

Det er verdt å dvele ved kor radikalt brotet med Bauhaus eigentleg var, for det vert ofte framstilt som ei vidareføring. Bauhaus hadde sin grunnkurs, sin \emph{Vorkurs}, der studenten skulle frigjere sansane sine og oppdage materialet sine eigne lover gjennom eksperiment — ein prosess som i botnen var kunstnarleg og personleg. Ulm sette i staden inn førelesingar i matematikk, logikk og vitskapsteori i dei same grunnåra. Der Bauhaus dyrka den skapande personlegdomen, ville Ulm helst gjere personlegdomen overflødig: ein god metode skulle gje gode resultat uavhengig av kven som brukte han, slik ei vitskapleg framgangsmåte gjev same svar i kven si hand som helst. Dette er ein nøkkel til heile Ulm-eksperimentet, og til kvifor det er så lærerikt. Skulen prøvde å flytte tyngdepunktet i faget frå \emph{utøvaren} til \emph{prosedyren} — frå den karismatiske formgjevaren med det gode auget, til ein etterprøvbar gang som kven som helst kunne følgje og kontrollere. Det er nett den flyttinga ein disiplin må gjere for å bli ein vitskap. Og det er nett den flyttinga som til slutt ikkje let seg fullføre.

\section{Modellen og marknaden}

Ulm-metoden vart ikkje berre undervist; han vart prøvd mot røynda, og det er her skulen sin ettermæle vart sikra. Gjennom eit langt samarbeid med elektronikkprodusenten Braun fekk Ulm syne kva metoden kunne gje. Samarbeidet var ikkje laust; ein eigen utviklingsgruppe ved skulen arbeidde direkte mot fabrikken, og resultatet vart ein heil produktfamilie som framleis vert rekna mellom det fremste industridesignet i hundreåret. Resultatet vart ein formverd av nøktern klårleik: radioar, platespelarar, barbermaskiner og kjøkenmaskiner reinsa for alt som ikkje tente bruken, ordna etter prinsipp ein kunne gjere greie for. Dieter Rams, som arbeidde tett på Ulm-miljøet, gav seinare denne haldninga sine mest kjende formuleringar, men røtene låg i Ulm sin tru på at god form var resultat av disiplin, ikkje av innfall. For ei kort tid såg det ut som om faget hadde funne noko: ein praksis der forma kunne grunngjevast, og der grunngjevinga gav former som heldt.

Og her må ein vere rettferdig mot Ulm, for det skulen oppnådde var verkeleg og varig. Dei former som kom ut av Braun-samarbeidet, er ikkje gått av mote slik strømlinja vart det; dei held framleis, fordi dei svara på noko meir varig enn ein sesong. Det er freistande å seie at metoden altså verka. Men sjå nøye på kva det var som gav formene deira styrke. Det var ikkje at metoden hadde funne ein objektiv godleik å optimere mot. Det var at miljøet delte ein \emph{smak} — ein streng, asketisk, nordisk-protestantisk smak for det reine og udekorerte — og at metoden gav denne smaken eit ordna, profesjonelt og overtydande uttrykk. Formene er gode etter den smaken, og smaken har vist seg robust. Men det var smaken som bar dei, ikkje metoden, og det er ein skilnad Ulm sjølv ikkje alltid heldt klår.

Men sjå nærare på kva som eigenleg var grunngjeve, og kva som ikkje var det. Metoden kunne ordne ein produksjonsprosess, analysere eit bruksmønster, optimere ein betjeningssekvens. Han kunne seie mykje om korleis ein knapp burde sitje for handa, og lite om kvifor heile apparatet burde finnast. Paul Betts har vist korleis Ulm-estetikken vart til ein nasjonal kulturell ressurs i Vest-Tyskland, eit prov på at den unge republikken hadde brote med fortida og funne ein rein, demokratisk, tiltrudd ting-kultur\autocite{betts2004}. Det er ein avslørande lagnad. Det Ulm meinte var resultatet av ein nøytral metode, vart oppfatta — og fungerte — som ein \emph{stil}: «god form», attkjenneleg, eksporterbar, med moralsk glans. Den same gliden vi såg i funksjonalismen, frå metode til smak utan at nokon merka overgangen, hende på nytt, midt i det mest metodebevisste miljøet faget hadde. Metoden disiplinerte detaljen. Han disiplinerte ikkje spørsmålet om kva som var verdt å lage, fordi til det fanst det ingen metode — berre den smaken metoden gav seg ut for å ha avskaffa.

\section{Kollapsen, dobbel}

I 1968 vart Ulm stengd. Den ytre årsaka var politisk og økonomisk: skulen var avhengig av offentleg løyving frå delstaten Baden-Württemberg, og eit fiendtleg politisk fleirtal kravde at skulen anten skulle leggjast inn under den tekniske høgskulen eller miste pengane. René Spitz har skrive den grundigaste skildringa av denne striden, og han syner at avviklinga var like mykje eit politisk attentat som ein fagleg dom: konservative krefter ville ikkje finansiere ein skule dei såg som ulydig, venstreorientert og økonomisk vanstyrt\autocite{spitz2002ulm}. Studentopprøret same året gjorde berre konflikten skarpare. Den institusjonelle Ulm døydde av pengemangel og politikk, slik mange gode institusjonar gjer, og det er freistande å stogge der: ein god idé drepen utanfrå.

Men det ville vere å lese halve historia. For samstundes med den ytre krisa hadde ei indre uro vakse fram, og ho råka sjølve metodeoptimismen. Fleire i miljøet, Maldonado mellom dei, hadde byrja å tvile på om den vitskaplege metoden verkeleg kunne bere det dei hadde lagt på han. Semiotikken og systemteorien gav ord og diagram, men dei gav ikkje formgjevaren det avgjerande: ein måte å avgjere kva form som var betre, når alle dei målbare krava var oppfylte og det framleis stod att eit val. Metoden kunne strukturere vegen fram til valet og så teie i sjølve valet. Han kunne fortelje deg korleis du skulle analysere problemet, men ikkje korleis du skulle vite at løysinga var rett, for det fanst ingen storleik å måle rettleiken mot.

Sjå kva dette tyder, for det er den faglege kjernen i kollapsen. Ein metode er ein gang ein følgjer; men ein gang fører berre fram til eit svar dersom det finst ein storleik gangen kan optimere — slik nedstigningsmetoden finn dalbotnen fordi det finst ei høgd å minke. Ulm hadde gangen, men ikkje storleiken. Det fanst ingen veldefinert «godleik» ved ei form som metoden kunne søkje mot, og difor enda kvar metodisk analyse i det same opne valet ho byrja med, berre betre opplyst. Studentane lærte å kartleggje eit problem grundig, og stod så att i nett det skjønnet metoden skulle ha avløyst. Det er ikkje at metodane var dårlege; det er at dei mangla ein funksjon å minimere, og ein metode utan ein storleik å optimere er ikkje ein metode i streng forstand, men ei sjekkliste. Den unge metoderørsla utanfor Ulm — folk som John Christopher Jones, som same året samla heile feltet av designmetodar i ei stor bok — skulle innan få år støyte på den same veggen, og fleire av grunnleggjarane kom til å vende seg vekk frå sin eigen metodisme i avsky\autocite{jones1970methods}. Det er eit sjeldsynt og ærleg syn: ei rørsle som forlét sitt eige prosjekt fordi ho såg at det ikkje heldt. Ulm rasa altså to gonger på ein gong: institusjonen utanfrå, og trua på metoden innanfrå.

\section{Kvifor det er lærerikt at det var her det fall}

Det er freistande å sjå Ulm som ein martyr — eit fag som var på rett veg og som vart drepen før det kom fram. Men den lesinga trøystar meir enn ho lærer. Den hardare og sannare lesinga er denne: Ulm kom nærast eit fundament, og avstanden som stod att då skulen fall, var ikkje liten. Skulen hadde gjort alt rett som let seg gjere med dei reiskapane tida hadde. Det hadde henta inn vitskapane, bygt ein metode, prøvd han mot industrien, og synt at metoden kunne gje verdig form. Og likevel mangla det same ting til slutt som det mangla i starten: eit omgrep som kunne seie, presist, kva ei form skal svare på, korleis dei mange og motstridande krava heng saman, og kva som gjer ei løysing betre enn ei anna når alle krava er oppfylte. Metoden organiserte prosessen kring eit tomrom han ikkje kunne fylle.

Dette peikar framover, mot ein draum som tok form same tida, men i ein heilt annan tradisjon. Same tiåret som Ulm stengde, gav Herbert Simon ut tanken om at det kunne finnast ein eigen \emph{vitskap om det kunstige} — ein streng, formell vitskap om korleis ein utformar ting som skal tene eit føremål\autocite{simon1969sciences}. Simon meinte alvor med ordet vitskap: han ville ha ein teori om designprosessen som kunne formaliserast, undervisast og prøvast, på line med ein naturvitskap, og han såg for seg at datamaskinen skulle gjere det mogleg å rekne på det som før hadde vore skjønn. Ulm var det praktiske, institusjonelle forsøket på nett denne draumen, gjort før draumen fekk sitt klåraste teoretiske uttrykk; og Ulm fall før Simon rakk å skrive han ned. Det er ein bitter ironi i kronologien: det mest seriøse praktiske forsøket på ein designvitskap rasa same tiåret som den mest seriøse teoretiske formuleringa av han vart fødd, og dei to fekk aldri møtast. Det neste kapitlet i fagets vitskaplege ambisjon vart difor ikkje ført vidare frå Ulm, men teke opp på nytt, frå teorien si side, av folk som ikkje visste eller ikkje brydde seg om at det praktiske forsøket alt hadde gått tom. Det er eit mønster vi kjem til å sjå att, og det er kanskje det mest øydeleggjande av alle: faget byrjar grunnlaget på nytt, generasjon for generasjon, fordi det aldri samla det førre forsøket sine lærdomar til noko som kunne byggjast vidare på. Eit fag der kvar generasjon må byrje på botnen, akkumulerer ikkje; og det som ikkje akkumulerer, er per definisjon ikkje ein vitskap, kor vitskapleg det enn kler seg.

Lærdomen frå Ulm er difor ikkje at metode er forgjeves, og heller ikkje at politikk drep gode skular, sjølv om begge er sanne. Lærdomen er at faget her, i sitt beste og mest disiplinerte forsøk, kom heilt fram til den veggen alle dei andre forsøka òg møtte, og at veggen ikkje var mangel på vilje eller pengar eller intelligens. Veggen var fråvêret av ein teori om sjølve forma. Ulm hadde metode utan eit fundament for metoden, slik funksjonalismen hadde ein formel utan innhald. Skilnaden er at Ulm var ærleg nok, og strengt nok, til å la fråvêret bli synleg.

Det er denne ærlegdomen som gjer Ulm uunnverleg i ein kritisk historie, og som skil skulen frå alt som kom etter. Dei seinare rørslene i faget — den semantiske vendinga, design thinking — møtte den same veggen og svara med å gjere veggen til ein dyd: dei erklærte at design er ein eigen kunnskapsform som ikkje treng den slags fundament, at problema er for vrange til å ha rette svar, at skjønnet er sjølve poenget. Ulm gjorde det motsette. Ulm trudde at fundamentet fanst, leita etter det med all den strengskap tida åtte, og fann det ikkje — og let det stå udekt at det ikkje var funne. Det er skilnaden mellom eit fag som ikkje har eit fundament og veit det, og eit fag som ikkje har eit fundament og har lært seg å kalle mangelen sin natur. Ulm høyrer til den fyrste sorten, og er difor det næraste designhistoria kjem ein heiderleg fallitt: ein konkurs der bøkene vart lagde fram opne, slik at ettertida kan sjå nett kva som mangla i kassa.

\vignett

Medan Ulm freista å gjere formgjevinga til ein vitskap som kunne grunngje seg sjølv, gjekk det på den andre sida av Atlanteren føre seg ei heilt anna utvikling. Der vart formgjevinga òg disiplinert av eit krav — men kravet var ikkje sanning eller metode. Det var sal. Og den disiplinen synte seg langt meir verknadsfull, og langt meir øydeleggjande, enn nokon metode Ulm fekk prøvd. Lat oss sjå kva som hende då designet vart ein forretningsmodell.

\vidarelesing{\fullcite{maldonado1958}; \fullcite{lindinger1991ulm}; \fullcite{spitz2002ulm}; \fullcite{betts2004}; \fullcite{jones1970methods}; \fullcite{simon1969sciences}}
